\chapter{Literature Review and Study of Existing Systems}

\section{Study of Existing Systems}
To establish an empirical baseline for contemporary grievance redressal paradigms, four dominant grievance management architectures were analyzed:
\begin{enumerate}
    \item \textbf{Centralised Public Grievance Redress And Monitoring System (CPGRAMS):} Developed by the Department of Administrative Reforms and Public Grievances (DARPG), Government of India, CPGRAMS provides a centralized platform for citizens to lodge grievances against Union Ministries and statutory bodies. While functionally comprehensive, CPGRAMS relies heavily on multi-tiered manual forwarding of grievances across bureaucratic silos. The portal lacks automated text classification, does not offer vernacular voice dictation, and frequently experiences multi-week processing delays.
    \item \textbf{Reserve Bank of India - Complaint Management System (RBI CMS):} Launched under the Reserve Bank - Integrated Ombudsman Scheme (RB-IOS), CMS provides a digital interface for lodging complaints against scheduled commercial banks, NBFCs, and payment system operators. While it represents a major regulatory advancement, CMS operates primarily as an appellate forum: customers can only approach CMS after exhausting 30 days of grievance pendency at their respective banks. The user interface is predominantly English-oriented, requires extensive manual form entry, and provides limited real-time predictive transparency into resolution queues.
    \item \textbf{Commercial Scheduled Bank Portals (e.g., SBI CMS, HDFC Grievance Portal):} Tier-1 commercial banks maintain proprietary web portals for ticket submission. These portals typically employ static dropdown menus for category selection. When customers misclassify their dispute (e.g., confusing an erroneous UPI chargeback with an unauthorized net banking breach), tickets are bounced between internal departments, leading to repeated SLA breaches. Furthermore, plain-text descriptions routinely store unmasked customer account and card details.
    \item \textbf{Commercial Enterprise Helpdesks (e.g., Zendesk, Freshdesk):} Widely adopted across corporate enterprises, these systems offer sophisticated ticketing and macro-based automation. However, they lack native adherence to Indian regulatory frameworks (such as RBI-mandated SLA timelines, banking ombudsman escalation hierarchies, and India's DPDP Act 2023 compliance). Furthermore, enterprise licensing fees render them economically unviable for smaller regional rural banks and urban co-operative credit institutions.
\end{enumerate}

\section{Review of Research Papers and Technical Resources}
A comprehensive survey of academic literature and technical standards was undertaken across multiple core computer engineering domains:
\begin{itemize}
    \item \textbf{Text Classification and Natural Language Processing:} McCallum and Nigam \cite{mccallum1998comparison} established the mathematical efficacy of Multinomial Naive Bayes (MNB) over multi-variate Bernoulli models for text categorization, demonstrating superior performance when dealing with document representations weighted by word frequency (such as TF-IDF). Pedregosa et al. \cite{pedregosa2011scikit} validated the computational efficiency of MNB on sparse matrices, making it exceptionally well-suited for high-throughput, low-latency API categorization workloads without necessitating expensive GPU clusters.
    \item \textbf{Operational Queueing Theory and Capacity Planning:} Kleinrock \cite{kleinrock1975queueing} formulated the foundational mathematical models for multi-server queueing systems ($M/M/c$). Applying Erlang-C waiting probability equations enables the analytical calculation of expected queue dwell time ($W_q$) and the probability of waiting longer than a prescribed SLA deadline ($P(W > t)$). These mathematical models provide rigorous theoretical justification for dynamic staffing allocation in banking operations.
    \item \textbf{Speech Interfaces and Vernacular Inclusivity:} Shires and Wennborg \cite{w3c_webspeech_2012} specified the W3C Web Speech API, defining client-side speech recognition interfaces that interface with underlying OS or browser speech synthesis and recognition engines. Research into vernacular human-computer interaction emphasizes that voice-driven interfaces reduce cognitive load and form-filling errors for non-English speakers by over 60\%.
    \item \textbf{Data Privacy and Regulatory Compliance:} The Digital Personal Data Protection Act (DPDP Act 2023) \cite{dpdp_act_2023} mandates strict technical safeguards against the unauthorized retention and exposure of personal financial identifiers. Automated tokenization and regular expression redaction at the API ingestion boundary satisfy statutory privacy requirements and eliminate leakage in distributed log aggregators.
\end{itemize}

\section{Comparative Analysis}
Table~\ref{tab:literature_comparison} provides a detailed comparative matrix evaluating GrievanceIQ against existing platforms across seven critical architectural and functional parameters.

\begin{table}[h]
\centering
\caption{Comparative Matrix: GrievanceIQ vs Existing Grievance Systems}
\label{tab:literature_comparison}
\begin{tabular}{|p{3.2cm}|p{2.6cm}|p{2.6cm}|p{2.6cm}|p{3.0cm}|}
\hline
\textbf{Evaluation Parameter} & \textbf{CPGRAMS} & \textbf{RBI CMS} & \textbf{Zendesk / Helpdesks} & \textbf{GrievanceIQ (Proposed)} \\ \hline
\textbf{Automated NLP Triage} & Manual triage & Rule-based dropdown & Paid add-on ML & \textbf{Native MNB (89.2\% accuracy)} \\ \hline
\textbf{Vernacular Voice Input} & Absent & Absent & Third-party plugin & \textbf{Native Web Speech (EN/HI/MR)} \\ \hline
\textbf{Multilingual UI} & English / Hindi & English / Hindi & Configurable & \textbf{Complete 100\% Parity (EN/HI/MR)} \\ \hline
\textbf{DPDP 2023 PII Masking} & None (Plain text) & None (Plain text) & Regex redaction & \textbf{Ingestion-Level Regex Scrubbing} \\ \hline
\textbf{Statutory PDF + QR} & Basic PDF & Acknowledgement & PDF receipt & \textbf{ReportLab Dynamic QR Receipt} \\ \hline
\textbf{Predictive Capacity ($M/M/c$)} & None & None & Basic linear extrapolation & \textbf{Analytical $M/M/c$ + What-If Sim} \\ \hline
\textbf{Evidence OCR Verification} & Manual review & Manual review & Generic attachment & \textbf{Automated 4-Point Entity Scoring} \\ \hline
\end{tabular}
\end{table}

\section{Limitations of Existing Systems}
\begin{itemize}
    \item \textbf{High Operational Cost and Slow Throughput:} Manual sorting and routing of tickets consume between 4 to 18 business hours before an operational officer begins investigation.
    \item \textbf{Accessibility Barriers:} Existing forms fail to support spoken vernacular languages, alienating millions of rural consumers who conduct transactions via UPI and AePS.
    \item \textbf{Regulatory Non-Compliance:} Unredacted Aadhaar, Debit Card numbers, and PAN records in support tickets expose institutions to stringent penalties under the DPDP Act 2023.
    \item \textbf{Lack of Explainability in Decisions:} Automated closures often provide boilerplate responses with zero transparent reasoning traces, inciting repeated escalations and disputes.
\end{itemize}

\section{Identified Research and Engineering Gap}
While machine learning algorithms, speech recognition APIs, and queueing models exist as isolated theoretical disciplines, no unified, production-grade banking platform bridges the divide between \textit{autonomous multi-modal citizen ingestion} (voice dictation, trilingual parity, PII stripping) and \textit{actionable back-office operational intelligence} (queue capacity modeling, explainable evidence verification, dynamic SLA monitoring).

\section{Proposed Improvements in GrievanceIQ}
GrievanceIQ bridges this engineering gap by:
\begin{enumerate}
    \item Deploying an end-to-end trilingual text categorization model trained on 15,000 real-world banking grievance records.
    \item Implementing a zero-latency, buffer-isolated Web Speech voice recognition engine that prevents repeated line repetition.
    \item Enforcing zero-trust PII sanitization at the API boundary, guaranteeing statutory compliance with the DPDP Act 2023.
    \item Synthesizing verifiable, tamper-evident statutory PDF receipts containing live QR codes for instantaneous citizen tracking.
    \item Formulating real-time queueing models ($M/M/c$) that deliver predictive foresight into operational backlogs and staffing requirements.
\end{enumerate}
